\section*{Chapter \ref{ch:b2:limb-organogenesis} --- Organogenesis: Building the Tetrapod Limb}

\begin{solution}{exo:b2:limb-organogenesis:1}
Proximo-distal: the \omterm{def:b2:limb-organogenesis:bud}{apical ectodermal ridge} along the tip, FGFs.
Antero-posterior: the \omterm{def:b2:limb-organogenesis:bud}{zone of polarising activity} at the posterior
margin, \omterm{prop:b2:limb-organogenesis:zpa}{Sonic hedgehog}. Dorso-ventral: the dorsal ectoderm, Wnt7a.
\end{solution}

\begin{solution}{exo:b2:limb-organogenesis:2}
Stylopod: humerus, femur (Hoxa/d 9--10). Zeugopod: radius and ulna,
tibia and fibula (Hox 11). Autopod: wrist and hand, ankle and foot
(Hox 13).
\end{solution}

\begin{solution}{exo:b2:limb-organogenesis:3}
Limb field specified (Tbx5); bud grows out under the ridge's FGF;
the ZPA and dorsal ectoderm pattern it; cells leaving the \omterm{prop:b2:limb-organogenesis:aer}{progress
zone} condense into rays; cartilage differentiates; joints form
where cartilage is suppressed; blood vessels invade and bone
replaces cartilage; interdigital cells die; muscle precursors from
the somites migrate in, fuse and attach by tendons.
\end{solution}

\begin{solution}{exo:b2:limb-organogenesis:4}
(a) A stump with little or no humerus; (b) humerus, radius and ulna
but no hand; (c) a complete limb --- FGF replaces the ridge.
\end{solution}

\begin{solution}{exo:b2:limb-organogenesis:5}
$x = 40\ln 2 = \qty{28}{\micro m}$, $40\ln 5 = \qty{64}{\micro m}$,
$40\ln 20 = \qty{120}{\micro m}$: territories of 28, 36 and
\qty{56}{\micro m}.
\end{solution}

\begin{solution}{exo:b2:limb-organogenesis:6}
(a) 4-3-2-2-3-4, a full mirror set. (b) A weaker source gives a
lower peak, so only the low-threshold digit is added: 2-2-3-4. (c)
A short exposure likewise specifies only the most anterior extra
digit: 2-2-3-4 --- cells integrate concentration over time.
\end{solution}

\begin{solution}{exo:b2:limb-organogenesis:7}
$5000\times 2^{8} = 1.3\times 10^{6}$. Reaching $4\times 10^{6}$ needs
$\log_2 800 = 9.6$ doublings, so the cycle must average about
\qty{10}{h}, or the count must include the muscle precursors that
migrate in from the somites; cell death would make the discrepancy
worse, not better.
\end{solution}

\begin{solution}{exo:b2:limb-organogenesis:8}
Chick interdigital cells receive a BMP signal and die; duck
interdigital cells express a BMP inhibitor and survive, so the web
persists. The inhibitor on a bead saves the chick web: a webbed
foot. Death on schedule in a graft means the cells were already
instructed before removal --- the decision is made early and
executed later, autonomously.
\end{solution}

\begin{solution}{exo:b2:limb-organogenesis:9}
Without Hoxa13 and Hoxd13: no autopod --- the forelimb ends at the
wrist. Without Hoxa11 and Hoxd11: a much shortened radius and ulna
with a hand attached almost directly to the humerus. Each segment
needs its own Hox pair; the domains are not merely markers but
instructions.
\end{solution}

\begin{solution}{exo:b2:limb-organogenesis:10}
FGF from the ridge keeps Shh on in the ZPA; Shh, through a mesodermal
relay, keeps FGF on in the ridge. Each centre's signal therefore
certifies that the other is working, so a limb never grows without
being patterned or is patterned without growing, and a transient
loss of either is repaired by the other. If Shh could not maintain
FGF, the ridge would regress within a day and outgrowth would stop
at whatever segment had been reached: a truncated limb. Labour's
loop is ended by the event it drives --- delivery removes the
stimulus; the limb's loop is ended by differentiation --- the
mesoderm stops relaying when the limb is complete.
\end{solution}

\begin{solution}{exo:b2:limb-organogenesis:11}
$e^{4r} = 8$: $r = \ln 8/4 = \qty{0.52}{d^{-1}}$. Growth: a longer,
faster proliferation of the digit cartilage (a limb enhancer that
raises a growth gene's expression in the digits). Death: survival
of the interdigital tissue through expression of a BMP inhibitor,
which keeps the membrane.
\end{solution}

\begin{solution}{exo:b2:limb-organogenesis:12}
Signals are reused because a working receptor--transduction module
is cheap to redeploy: what changes between tissues is not the
signal but the set of genes its receiving cells are competent to
switch on, so the same concentration profile means motor neuron in
one place and digit 4 in another. A \omterm{def:b2:genome-diversification:mutation}{mutation} in such a gene
therefore affects many organs at once --- Shh \omterm{def:b2:genome-diversification:mutation}{mutations} give
holoprosencephaly and limb defects together --- unless it lies in a
tissue-specific enhancer, in which case it changes one organ only,
which is how most evolutionary change in these genes occurs.
\end{solution}

\begin{solution}{pb:b2:limb-organogenesis:1}
\textbf{1.} $96/12 = 8$ doublings: $2000\times 256 = 5.1\times 10^{5}$
cells.
\textbf{2.} $1000/256 \approx 4$: cell enlargement and cartilage
matrix.
\textbf{3.} Length $\times 10$; a cylinder of constant radius would
give $\times 10$ in volume; the extra hundredfold means the radius
also grew about tenfold --- the bud widened into a paddle as it
lengthened.
\textbf{4.} At 3 days $300/400 = \qty{75}{\%}$; at 7 days $300/4000 =
\qty{7.5}{\%}$.
\textbf{5.} Only the tip is under the signal; as the tip advances,
cells left behind leave the zone in the order they were made, so
proximal structures are specified first and distal ones last.
\textbf{6.} The zeugopod (radius and ulna), specified between 3.5 and
4 days.
\textbf{7.} Stylopod by 3.5 days, zeugopod by 4.0 days, autopod by 4.5
days (digit tips later).
\textbf{8.} Half a day: one cell cycle.
\textbf{9.} Each cycle spent in the \omterm{prop:b2:limb-organogenesis:aer}{progress zone} advances the Hox
code (9--10, then 11, then 13); cells leaving after one, two or
three cycles carry the code of the stylopod, zeugopod or autopod.
\textbf{10.} A complete wing: FGF is what the ridge provided.
\textbf{11.} A second outgrowth from the dorsal surface: duplicated
distal structures, a supernumerary limb.
\textbf{12.} The mesoderm stops relaying Shh to the ridge, Shh
expression ends, the ridge regresses, and cells differentiate
instead of dividing: the loop that sustained growth is broken.
\textbf{13.} $k = \ln 2/1740 = \qty{4.0e-4}{s^{-1}}$; $\lambda =
\sqrt{10^{-12}/4\times 10^{-4}} = \qty{50}{\micro m}$.
\textbf{14.} $x = 50\ln(1/0.6) = 26$, $50\ln 4 = 69$, $50\ln 12.5 =
\qty{126}{\micro m}$: territories of 26, 43 and \qty{57}{\micro m}.
\textbf{15.} $600 - 126 = \qty{474}{\micro m}$.
\textbf{16.} Every boundary moves out by $50\ln 2 = \qty{35}{\micro
m}$: 61, 104, 161; the digit-2 territory grows from 57 to
\qty{92}{\micro m}, enough for an extra anterior digit 2.
\textbf{17.} Each set needs \qty{126}{\micro m}: at least
\qty{252}{\micro m}; the \qty{600}{\micro m} bud has room, with
\qty{348}{\micro m} of digit-free tissue in the middle.
\textbf{18.} At \qty{200}{\micro m} the two gradients overlap and
add: the middle cells see enough for digit 3 and the digit-2
territories vanish --- 4-3-4 or 4-3-3-4, fewer digits than a full
mirror set.
\textbf{19.} A quarter of the cells gives $c_0/4$: no cell reaches
the threshold of digit 4 or 3 ($0.25 c_0$ at the margin itself), and
the digit-2 threshold is passed out to $50\ln(0.25/0.08) =
\qty{57}{\micro m}$: one extra digit 2.
\textbf{20.} $3\times 300\times 300\times 200 = 5.4\times 10^{7}$
\unit{\micro m^3}: $5.4\times 10^{4}$ cells; over \qty{86400}{s},
$0.6$ cells per second.
\textbf{21.} BMP bead in the duck web: the web dies and the toes
separate; inhibitor in the chick web: the web survives, a webbed
foot.
\textbf{22.} $e^{0.5\times 4} = e^{2} = 7.4$. The second change is
survival of the interdigital membrane by a BMP inhibitor.
\textbf{23.} The late phase of Hox 13 expression that specifies an
autopod, and the persistence of the ridge: the fish's fold makes fin
rays instead. \emph{Tiktaalik}'s fin contains wrist-like bones --- a
fin beginning to be a limb.
\textbf{24.} $5.4\times 10^{4}/1000 = 54$: $\log_2 54 = 5.8$ doublings,
about \qty{69}{h}, three days to make what one day removes.
\textbf{25.} Eight doublings; digit boundaries at 26, 69 and
\qty{126}{\micro m}; a mirror duplication needs at least
\qty{252}{\micro m}.
\end{solution}
